\section{Conclusion and Future Work}\label{sec:conclusion}

This paper presented a blockchain-based architectural model for disease early warning and surveillance, implemented as a functional prototype on the Ethereum platform. The model integrates four core surveillance processes---hierarchical disease case reporting, laboratory result reporting, automated report completeness verification, and multi-dimensional epidemiological data analysis---within a decentralized, immutable, and traceable smart contract architecture.

The three-dimensional evaluation validated the model against operational requirements. Performance evaluation demonstrated stable block generation times of 12--14 seconds under Proof of Authority consensus, satisfying the temporal requirements of surveillance reporting workflows. Usability evaluation via the standardized SUMI instrument ($n=26$ domain professionals) yielded scores of 60.46--68.96 across all six scales, significantly exceeding the population mean and confirming practical usability for epidemiologists, public health workers, and laboratory technicians. Security evaluation through threat modeling demonstrated inherent resistance to man-in-the-middle, denial-of-service, and traffic analysis attacks through blockchain's native cryptographic and decentralization properties.

The proposed model addresses the identified limitations of existing disease surveillance systems---data centralization, security vulnerabilities, traceability gaps, manual completeness verification, and inter-organizational interoperability barriers---through the systematic application of distributed ledger technology to the public health domain.

Future research directions include:
\begin{itemize}
\item Extension of the model to encompass public health emergency response and recovery processes downstream of surveillance.
\item Implementation of incidence proportion and prevalence rate computation functions within the smart contract architecture.
\item Upgrade of the prototype from test network deployment to a production-grade operational system.
\item Evaluation of alternative blockchain platforms beyond the three compared in this study.
\item Integration of machine learning and artificial intelligence modules with the blockchain architecture for predictive surveillance capabilities.
\item Development of a cross-platform mobile application to improve field accessibility.
\item Comparative evaluation against similar blockchain-based health information systems.
\end{itemize}

%% ===== BACK MATTER =====

\section*{Author Contributions}
\textbf{Teklhiwot Mekonen:} Conceptualization; Methodology; Software; Validation; Formal analysis; Investigation; Data curation; Writing -- original draft; Visualization.

\textbf{Esubalew Alemneh:} Supervision; Writing -- review \& editing.

\section*{Acknowledgments}
The authors thank the staff of the Amhara Public Health Institute for their participation in the blockchain suitability assessment survey and the SUMI usability evaluation. The authors also acknowledge the Human Computer Interaction Research Group (HCIRG), University College Cork, Ireland, for provision of the SUMI academic license.

\section*{Financial Disclosure}
This research was conducted as part of a Master of Science thesis at Bahir Dar University, Bahir Dar Institute of Technology, Faculty of Computing. No specific external grant from any funding agency in the public, commercial, or not-for-profit sectors was received for this research.

\section*{Conflict of Interest}
The authors declare no conflicts of interest.

\section*{Data Availability Statement}
The smart contract source code for the prototype is available in the repository accompanying this manuscript. The SUMI questionnaire instrument is proprietary and available under academic license from the Human Computer Interaction Research Group, University College Cork, Ireland. Raw SUMI response data are not publicly available due to participant confidentiality but are available from the corresponding author upon reasonable request.

\section*{Ethics Statement}
This study involved a usability evaluation with 26 human participants (13 field epidemiologists, 7 public health workers, and 6 laboratory technicians) recruited from the Amhara Public Health Institute, Bahir Dar, Ethiopia. Participants were briefed on the research objectives and provided voluntary agreement to participate. The study involved interaction with a software prototype for 15--20 minutes followed by completion of the standardized Software Usability Measurement Inventory (SUMI) questionnaire. The study posed no more than minimal risk to participants: no personal health data were collected, no clinical interventions were performed, and all questionnaire responses were anonymous.


